% !TEX program = xelatex
% =============================================================================
% cn-legal-opinion · 补充法律意见书样张（问询回复式、信函版式）
%
% 版式对标通力/君合的信函式补充意见书：无封面，首页信头（律所名＋地址块，
% 不加横线）后直接题名行文；反馈意见原文整段加粗、「回复：」引出作答
% （通力 2018 宁德时代补充意见书体例）；页脚左下斜体案号（通力式）。
% 与 main.tex 共用 opinion.sty。
%
% 编译：bash scripts/compile.sh templates/cn-legal-opinion/supplemental.tex --preview
% =============================================================================
\documentclass[zihao=-4]{ctexart}
\usepackage{opinion}

% --- 文档元信息 --------------------------------------------------------------
\renewcommand{\opFirm}{北京市示例律师事务所}
\renewcommand{\opFirmEn}{EXAMPLE \& PARTNERS}
\renewcommand{\opClient}{示例科技股份有限公司}
\renewcommand{\opMatter}{首次公开发行股票并在科创板上市}
\renewcommand{\opConn}{的}   % 与 main.tex（金杜式「的」）保持一致
\renewcommand{\opDocType}{补充法律意见书（一）}
\renewcommand{\opRefNo}{示例（证）字〔2026〕第0001-1号}
\renewcommand{\opFascicle}{}                    % 信函式：纯页码
\renewcommand{\opMatterNo}{2026001/IPO/supp/D1}   % 左下斜体案号（通力式）

\begin{document}

% --- 信头＋题名（通力/君合式信函首页，无封面）--------------------------------
\opLetterhead
\opLetterTitle

\opTo{示例科技股份有限公司}

根据示例科技股份有限公司（以下简称"发行人"）的委托，本所指派王××律师、
张××律师（以下合称"本所律师"）担任发行人首次公开发行股票并在科创板上市
（以下简称"本次发行上市"）的专项法律顾问，已就本次发行上市出具了示例
（证）字〔2026〕第0001号《北京市示例律师事务所关于示例科技股份有限公司首
次公开发行股票并在科创板上市的法律意见书》（以下简称"《法律意见书》"）。

现根据上海证券交易所于2026年7月10日出具的《关于示例科技股份有限公司首次
公开发行股票并在科创板上市申请文件的审核问询函》（以下简称"《问询函》"）
的要求，本所律师对《问询函》所涉法律事项进行了进一步核查验证，特出具本补
充法律意见书。《法律意见书》中所述及的本所及本所律师的声明事项以及相关词
语定义同样适用于本补充法律意见书，但本补充法律意见书另作定义的除外。本补
充法律意见书构成《法律意见书》的补充。

% --- 问询回复（通力/中伦式：问题原文加粗，「回复：」引出作答）---------------
\opFeedback{1}{关于实际控制人认定。招股说明书披露，发行人实际控制人为林×
×，直接及间接合计控制发行人58.60\%的表决权。请发行人说明：(1)报告期内公
司股权结构的变动情况及实际控制人认定依据；(2)其他持股5\%以上股东与实际控
制人之间是否存在一致行动关系或者其他利益安排。请保荐机构、发行人律师核查
并发表明确意见。}

\opReply

\subsection{报告期内公司股权结构的变动情况及实际控制人认定依据}

经本所律师核查发行人的工商登记档案、历次股东会决议及股东出资凭证，报告期
内发行人共发生两次增资、一次股权转让，各次变动均已履行法定程序并办理变更
登记。截至本补充法律意见书出具之日，林××直接持有发行人42.35\%的股份，并
通过担任示例合伙企业（有限合伙）执行事务合伙人间接控制发行人16.25\%的表决
权，合计控制58.60\%的表决权，能够对股东会决议产生重大影响，且报告期内一直
担任发行人董事长、总经理，主导发行人的重大经营决策。据此，本所认为，将林
××认定为发行人的实际控制人依据充分，报告期内发行人实际控制人未发生变更。

\subsection{其他持股5\%以上股东与实际控制人之间的关系}

经本所律师核查其他持股5\%以上股东的调查表、出资来源说明及其与实际控制人签
署的确认函，该等股东与林××之间不存在一致行动协议、委托表决安排或者其他
利益安排。

\hl{综上，本所认为，发行人实际控制人认定准确，其他持股5\%以上股东与实际控
制人之间不存在一致行动关系或者其他利益安排。}

\opFeedback{2}{关于关联交易。招股说明书披露，报告期内发行人向关联方示例智
算科技有限公司采购算力服务。请发行人说明该等关联交易的定价公允性及必要性，
是否存在对发行人利益的损害。请发行人律师核查并发表明确意见。}

\opReply

经本所律师核查该等关联交易的合同文本、定价依据及同期市场可比价格，报告期
内发行人向示例智算科技有限公司采购算力服务的金额分别为320.00万元、
560.00万元和480.00万元，占各期营业成本的比例均低于3\%；交易价格参照同期
市场公开报价确定，与独立第三方报价不存在显著差异；该等交易已按《公司章
程》及《关联交易管理制度》的规定履行了审议程序，关联股东回避表决。

\hl{综上，本所认为，上述关联交易具有商业合理性，定价公允，已履行必要的审
议程序，不存在损害发行人及其他股东利益的情形。}

\vspace{0.5\baselineskip}
本补充法律意见书正本肆份，无副本，经本所盖章并经办律师签字后生效。

% --- 签署页 ------------------------------------------------------------------
\opSigPage{%
  \opsigblock{负责人：}{\opsigrow{林××}{}}
  \vspace{0.6\baselineskip}
  \opsigblock{经办律师：}{\opsigrow{王××}{张××}}}

\end{document}
